\chapter{CONSIDERAÇÕES FINAIS} \label{cap:conclusoes}

\section{Conclusões}

Este trabalho analisou como o mercado de trabalho brasileiro tem se transformado com o deslocamento setorial em direção ao setor terciário. A revisão teórica mostrou que, de acordo com a Escola Austríaca e Schumpeter, a inovação e o empreendedorismo impulsionam essas mudanças, criando lucros acima da média e reorganizando estruturas produtivas \cite{iorio2013, silva2025}.

Conceitos básicos (fatores de produção, agentes econômicos) e a classificação por CNAE permitiram contextualizar esses fenômenos no Brasil. A transição de trabalhadores do campo e da indústria para os serviços foi evidenciada por dados do IBGE e da literatura, que apontam que mais de dois terços da força de trabalho atua no terciário \cite{sescap2017, politize2024a}.

O estudo de caso da M. Dias Branco ilustrou como uma grande indústria gera emprego direto significativo, mas também demanda serviços de apoio e induz o emprego indireto. Por fim, conclui-se que, embora a indústria gere empregos diretos valiosos, o setor terciário ampliado (considerando comércio e serviços correlatos) cria o maior número total de empregos no nosso país \cite{sescap2017, dgabc2002}. Essa constatação reforça a importância de políticas públicas que estimulem não só a competitividade industrial, mas também a qualificação e a expansão do setor de serviços, equilibrando o crescimento econômico com a sustentabilidade social.


